\naslov{Analiza besedil in vpetja}

Širše obravnavamo poljubne delno strukturirane podatke $O$, zraven tega pa imamo
še slovar $E$.
Najenostavnejša ideja je, da poiščemo število pojavitev vsake besede v $O$, s
čimer dobimo vektor velikosti $\abs{E}$ s komponentami v $\N_0$.
Če želimo, lahko te podatke še transformiramo, v primeru besedil recimo
\begin{itemize}
\item Boolova frekvenca: $1$, če se beseda pojavi, sicer $0$
\item Normalizirana frekvenca: delimo s številom besed v besedilu
\item Logaritmična frekvenca: $\log(1 + f)$
\end{itemize}
Problem pri uporabi frekvenc za besedila je, da imamo veliko \pojem{praznih
  besed}, tj.~besed, ki si pojavijo velikokrat v vsakem besedilu.
Ena možna rešitev je, da praznih besed ne damo v slovar, druga pa, da pogledamo
frekvenco v dokumentih $\text{DF}$, tj.~število dokumentov, kjer se ta
beseda pojavi.
Potem definiramo mero
\[
  \text{TFIDF} = \text{TF} \cdot \text{IDF}
\]
za obratno frekvenco v dokumentih
\[
  \text{IDF} = \log \frac{\text{število dokumentov}}{\text{DF}}.
\]

\vprasanje{Opiši obravnavo vpetij s frekvencami. Kaj je TFIDF?}

Pri besedilih je preprocesiranja še veliko več.
Prvo je lematizacija, kjer pretvorimo vse besede v osnovno obliko (odstranimo
sklanjanje, množino, \ldots).
Če lematizacije ne moremo uporabiti, jo lahko zamenjamo s krnjenjem, kjer besedo
pretvorimo v njen koren.

\vprasanje{Kaj je lematizacija in kaj krnjenje?}

V primeru besedil slovar dobimo enostavno tako, da zberemo vse besede v
dokumentih.
Včasih dodamo tudi pogoste besedne zveze; temu potem rečemo \pojem{$n$-gram}.
V primeru slik so besede konvolucijski filtri; te lahko poiščemo na roko ali
uporabimo neko iskanje pogostih vzorcev.

Težava pri tej obravnavi bag-of-words je, da ignorira pomen besed.
Rešitev temu so \pojem{vpetja}, tj.~preslikave $E \to \R^n$, kjer pričakujemo,
da se bodo podobne besede slikale v podobne vektorje.
Pri tem naredimo \pojem{porazdelitveno predpostavko}: besede s podobnih pomenom
uporabimo v istem kontekstu.
Tu je kontekst mišljen kot nekaj besed levo in desno od trenutno obravnavane.

Za vhod v nevronske mreže besede kodiramo kot vektor dolžine $\abs{E}$ z enico
na indeksu, ki ustreza tej besedi v slovarju.
Model CBOW (\textit{continuous bag-of-words}) je preprosta nevronska mreža s
tremi plastmi.
Vhodna raven ima $2 k \abs{E}$ nevronov, ki kodirajo besede v kontekstu ($k$ je
širina konteksta).
Skrita raven zakodira predstavitveni vektor.
Ima $m$ nevronov, kjer je $m$ dimenzija vpetja.
Izhodna raven pa kodira opazovano besedo, spet ima $\abs{E}$ nevronov.
Pri tem naredimo eno poenostavitev; nevroni v vhodni plasti so razdeljeni v
$2k$ skupin, uteži na sinapsah so enake med skupinami (vsaka skupina kodira eno
besedo).
Ideja je, da to mrežo uporabimo tako, da poskuša napovedati srednjo besedo iz
besed konteksta.

\vprasanje{Opiši model CBOW.}

Lahko napovedujemo tudi v drugo smer; iz besede ustvarimo kontekst.
To naredimo z modelom Skip-Gram, ki je CBOW v nasprotni kategoriji.
Vektor stanj nevronov v skriti plasti nam bo tako dal vpetje vhodne besede.
Ciljna funkcija je sedaj negativni logaritem verjetja konteksta; pri izpeljavi
pa naredimo hudo napako, da so besede v kontekstu neodvisne.

\vprasanje{Opiši model Skip-Gram.}

% LocalWords:  lematizacija lematizacije konvolucijski bag-of-words CBOW
% LocalWords:  continuous Skip-Gram
